%! Function Operations — Unit 1, Lesson 1.2 — Student Packet Cover Sheet
% The ONE place \namedateperiod belongs (Namestrip).
\documentclass[10pt]{article}
\usepackage{precalculus-article}
\usepackage{precalculus-boxes}
\usepackage{ltablex}
\keepXColumns

\begin{document}

% Full-bleed plum banner
\begin{tikzpicture}[remember picture, overlay]
  \fill[plum] (current page.north west)
    rectangle ([yshift=-0.9in]current page.north east);
\end{tikzpicture}
\vspace{-0.8in}

\begin{center}
  {\color{white}\textbf{\LARGE Precalculus}}\\[4pt]
  {\color{lilacmid}\large\textbf{Unit 1: Functions and Change}}\\[6pt]
  {\color{white}\textbf{\large Lesson 1.2 \quad Function Operations}}
\end{center}

\vspace{0.22in}
\namedateperiod
\vspace{0.18in}

\begin{learningtargetbox}
\small
\begin{enumerate}[leftmargin=1.5em, itemsep=5pt]
  \item I can add, subtract, multiply, and divide two functions, and explain that
        each operation is arithmetic done to the two \textit{outputs}.
  \item I can evaluate a combination such as $(f+g)(2)$ at a point from a table
        or from formulas, and explain why $(f-g)$ and $(g-f)$ are different.
  \item I can build the formula for a combination and simplify it, keeping each
        function's whole rule inside parentheses.
  \item I can state the domain of a combination from the \textit{original}
        functions --- the overlap of the two domains, with every input that makes
        the bottom function zero thrown out of a quotient.
\end{enumerate}
\end{learningtargetbox}

\vspace{0.14in}

\begin{tocbox}
\small
\renewcommand{\arraystretch}{1.5}
\begin{tabularx}{\linewidth}{@{} c l X r @{}}
\rowcolor{lilac}
\textbf{\#} & \textbf{Component} & \textbf{Description} & \textbf{Score} \\
\midrule
1 & Warm-Up      & Spiral review of prerequisite skills            & \blank{1.2cm} \\
2 & Guided Notes & Combining functions at a point and as a formula, and the domain of the result & \blank{1.2cm} \\
\rowcolor{goldbg}
3 & AP Practice  & \textbf{Extra credit} --- AP-style multiple choice and free response & $+$\,\blank{1.0cm} \\
4 & Homework     & Graded practice in new contexts \quad Due: \blank{2.2cm} & \blank{1.2cm} \\
\midrule
  & \multicolumn{2}{r}{\textbf{Total} \quad {\footnotesize (1, 2, and 4, plus any extra credit)}} & \blank{1.2cm} \\
\end{tabularx}
\end{tocbox}

\vspace{0.14in}

\begin{remindbox}
\small
Two functions that take the same input can be combined into a \textbf{third}
function four ways --- add, subtract, multiply, divide --- and every one of them
is just arithmetic done to the two \textbf{outputs}. Order matters: $f - g$ and
$g - f$ are opposites, not the same function. And the new function is choosier
about its inputs than either parent: it accepts only what \textbf{both} parents
accept, and a quotient additionally refuses every input that makes the bottom
function \textbf{zero}. That refusal is permanent --- simplifying the formula
afterward can hide a forbidden input, but it can never make it legal.
\end{remindbox}

\end{document}
